\begin{tikzpicture}[
  font=\rmfamily\footnotesize,
  >=Latex,
  system/.style={
    draw=blue!60!black,
    fill=blue!7,
    rounded corners=2pt,
    align=center,
    text width=31mm,
    minimum height=15mm,
    inner xsep=2.5mm,
    inner ysep=2.5mm
  },
  attacker/.style={
    draw=red!65!black,
    fill=red!7,
    rounded corners=2pt,
    align=center,
    text width=38mm,
    minimum height=15mm,
    inner xsep=2.5mm,
    inner ysep=2.5mm
  },
  transition/.style={
    attacker,
    text width=42mm,
    fill=orange!9,
    draw=orange!75!black
  },
  asset/.style={
    draw=green!55!black,
    fill=green!9,
    rounded corners=2pt,
    align=left,
    text width=42mm,
    minimum height=14mm,
    inner xsep=3mm,
    inner ysep=2.5mm
  },
  corpus/.style={
    draw=blue!60!black,
    fill=blue!7,
    rounded corners=2pt,
    align=center,
    text width=35mm,
    minimum height=14mm,
    inner xsep=2.5mm,
    inner ysep=2.5mm
  },
  flow/.style={->,thick,draw=black!75},
  attackflow/.style={->,thick,draw=red!65!black},
  conditionflow/.style={->,thick,densely dashed,draw=orange!80!black}
]
  % Principal system path: one direction and no edge labels.
  \node[system] (interface) at (-8.0,0) {
    \textbf{Query interface}\\
    Current role and mode
  };
  \node[system] (retrieval) at (-4.0,0) {
    \textbf{Index and retrieval}\\
    Search and permitted relations
  };
  \node[system] (context) at (0,0) {
    \textbf{Model-visible context}\\
    Projected fields and retained state
  };
  \node[system] (generation) at (4.0,0) {
    \textbf{Generator}\\
    Raw answer
  };
  \node[system] (delivery) at (8.0,0) {
    \textbf{Delivery controls}\\
    Guard action and returned answer
  };

  \draw[flow] (interface) -- (retrieval);
  \draw[flow] (retrieval) -- (context);
  \draw[flow] (context) -- (generation);
  \draw[flow] (generation) -- (delivery);

  % Corpus input is kept within the system boundary.
  \node[corpus,below=13mm of retrieval] (corpus) {
    \textbf{Structured corpus}\\
    Indexed entities and synthetic rows
  };
  \draw[flow] (corpus) -- (retrieval);

  \node[
    draw=gray!65,
    densely dotted,
    rounded corners=2mm,
    fit=(interface)(retrieval)(context)(generation)(delivery)(corpus),
    inner xsep=4mm,
    inner ysep=4mm,
    label={[font=\rmfamily\scriptsize,text=gray!70!black]below right:SYSTEM UNDER TEST}
  ] (systemboundary) {};

  % Attacker capabilities are stated in boxes so arrows remain unobstructed.
  \node[attacker,above=15mm of interface] (chat) {
    \textbf{Standard chat attacker}\\
    Public or internal role\\
    Natural-language requests: A01--A05, A08
  };
  \node[transition,above=15mm of context] (rolechange) {
    \textbf{Authority-transition condition}\\
    Experiment-controlled high-to-low change: A03
  };
  \node[attacker,below=15mm of corpus] (poison) {
    \textbf{Poisoned-corpus attacker}\\
    Attacker-controlled synthetic public entity: A06/A07 only
  };

  \draw[attackflow] (chat) -- (interface);
  \draw[conditionflow] (rolechange) -- (context);
  \draw[attackflow] (poison) -- (corpus);

  % Security and utility outcomes are distinct evaluation targets.
  \node[asset,right=15mm of delivery,yshift=20mm] (conf) {
    \textbf{Confidentiality assets}\\
    Fields, rows, relation edges, protected indexed-record existence, and conversation state
  };
  \node[asset,below=6mm of conf] (integrity) {
    \textbf{Integrity asset}\\
    Trusted-instruction boundary and canary-free delivery
  };
  \node[asset,below=6mm of integrity] (utility) {
    \textbf{Authorised task}\\
    Protected-role expected answer
  };

  \draw[flow] (delivery.east) -- ++(5mm,0) |- (conf.west);
  \draw[flow] (delivery.east) -- (integrity.west);
  \draw[flow] (delivery.east) -- ++(5mm,0) |- (utility.west);

  \node[
    draw=green!45!black,
    densely dotted,
    rounded corners=2mm,
    fit=(conf)(integrity)(utility),
    inner sep=3mm,
    label={[font=\rmfamily\scriptsize,text=green!45!black]above left:EVALUATED OUTCOMES}
  ] {};
\end{tikzpicture}
